\chapter{Cow's Milk Allergy}
\index{Cow's Milk Allergy}

\section*{Definition and Overview}
Cow's milk allergy is an immune reaction to proteins in cow's milk, most commonly in babies and young children. It is different from lactose intolerance, which is a digestive problem rather than an allergy. Symptoms can affect the skin, stomach and bowel, and, in severe cases, the airways and circulation. Management involves strictly avoiding cow's milk and using suitable alternatives, and most children outgrow the allergy by school age. Because cow's milk is common in foods, diagnosis and dietary advice are best provided with specialist and dietitian support.

\section*{Epidemiology}
Cow's milk allergy is one of the most common food allergies in early childhood, affecting around 2--3\% of infants and young children. It usually appears in the first year of life, often when formula is introduced, and in breastfed babies whose mothers consume dairy. Most children outgrow cow's milk allergy: around half by age five and the majority by adolescence, although those with immediate, severe reactions are less likely to outgrow it. In many countries, as elsewhere, the incidence appears to have risen over recent decades.

\section*{Aetiology and Pathophysiology}
Cow's milk contains several proteins, chiefly casein and whey proteins, that can trigger allergic reactions.

\begin{itemize}
    \item \textbf{IgE-mediated allergy:} the immune system produces antibodies (IgE) to milk proteins, causing rapid reactions within minutes to two hours
    \item \textbf{Non-IgE-mediated allergy:} delayed reactions involving other immune cells, with symptoms appearing hours to days later
    \item \textbf{Casein and whey proteins:} the main allergens, present in most dairy products
    \item \textbf{Sensitisation:} exposure to milk proteins triggers the allergic response
    \item \textbf{Anaphylaxis:} a severe, rapid allergic reaction affecting breathing and blood pressure, which can be life-threatening
    \item \textbf{Eosinophilic disorders:} rare inflammatory conditions of the gut related to food allergy
\end{itemize}

\section*{Clinical Features}
Symptoms depend on whether the reaction is immediate or delayed.

\begin{itemize}
    \item \textbf{Immediate reactions:} hives, swelling of the lips or face, vomiting, and, in severe cases, difficulty breathing or collapse
    \item \textbf{Delayed reactions:} eczema, reflux, colic-like symptoms, vomiting, diarrhoea, blood in the stool, poor weight gain, and constipation
    \item \textbf{Respiratory symptoms:} wheeze or a runny nose in some children
    \item \textbf{Symptoms within minutes to two hours of ingestion:} typical of IgE-mediated allergy
    \item \textbf{Recovery with avoidance:} symptoms settle when milk is removed from the diet
\end{itemize}

\section*{Differential Diagnosis}
Many conditions mimic cow's milk allergy, especially in babies.

\begin{itemize}
    \item Lactose intolerance, causing wind, diarrhoea, and bloating but not skin or respiratory symptoms
    \item Gastro-oesophageal reflux and colic, common in babies for other reasons
    \item Eczema not related to food allergy
    \item Gastroenteritis and other gut infections
    \item Other food allergies, such as egg or soy allergy
    \item Coeliac disease or other malabsorption conditions
    \item Constipation from other causes
\end{itemize}

\section*{When to Seek Medical Care}
If you suspect your baby or child has a reaction to cow's milk, see a GP or paediatrician for assessment rather than removing milk without advice, because dietary changes need to be safe and nutritionally complete.

\textbf{Contact your local emergency services for:}
\begin{itemize}
    \item Difficulty breathing, wheezing, or swelling of the face, lips, or throat
    \item A severe allergic reaction (anaphylaxis): collapse, confusion, or vomiting with breathing difficulty
    \item Signs of severe dehydration from vomiting or diarrhoea
    \item A child who is floppy, very drowsy, or difficult to wake
\end{itemize}

\section*{Diagnosis and Investigations}
Diagnosis combines history, examination, and specific tests.

\begin{itemize}
    \item \textbf{History:} the timing of symptoms relative to milk exposure, and their nature
    \item \textbf{Examination:} assessment of growth, skin, and general health
    \item \textbf{Skin prick testing:} a small drop of milk protein placed on the skin; a wheal indicates sensitisation
    \item \textbf{Blood test:} specific IgE (RAST) to milk proteins
    \item \textbf{Trial of avoidance:} removing milk for two to four weeks to see if symptoms settle
    \item \textbf{Food challenge:} supervised reintroduction in hospital or clinic for uncertain cases
    \item \textbf{Referral:} to an allergist or paediatrician for complex cases
\end{itemize}

\section*{Management}
Management is strict avoidance, with careful attention to nutrition.

\begin{itemize}
    \item \textbf{Breastfeeding:} mothers may continue breastfeeding and avoid cow's milk themselves, with calcium and vitamin D supplements
    \item \textbf{Formula:} extensively hydrolysed or amino acid formulas for babies who cannot tolerate cow's milk formula
    \item \textbf{Soy formula:} an option for some children, depending on age and tolerance
    \item \textbf{Food avoidance:} checking labels for milk, casein, whey, and other milk-derived ingredients
    \item \textbf{Emergency plan:} antihistamines for mild reactions and an adrenaline autoinjector for children at risk of anaphylaxis
    \item \textbf{Dietitian advice:} to ensure calcium, protein, and energy needs are met
    \item \textbf{Regular review:} allergy tests may be repeated, and supervised challenges arranged to test for outgrowing
    \item \textbf{Education:} family, childcare, and school awareness of the allergy and emergency plan
\end{itemize}

\section*{Complications}
\begin{itemize}
    \item Anaphylaxis, a life-threatening emergency
    \item Poor growth and nutrition if milk is replaced without appropriate alternatives
    \item Calcium and vitamin D deficiency with low bone density in later life
    \item Eczema that is difficult to control
    \item Iron deficiency, as cow's milk avoidance can affect diet
    \item Anxiety in the child and family about accidental exposure
    \item Missed diagnosis when symptoms are attributed to other conditions
\end{itemize}

\section*{Prognosis and Outcomes}
The majority of children with cow's milk allergy outgrow it, most by school age, and non-IgE-mediated allergy has the best outlook. Children with immediate, severe reactions are more likely to remain allergic into adolescence, and re-testing is used to decide when a supervised food challenge is safe. With good dietary management, growth and development are normal. A small proportion of children with multiple allergies or severe eczema have a longer course and need continued specialist care.

\section*{Prevention and Self-Care}
\begin{itemize}
    \item Introduce common foods, including those containing milk, by around six months (not before four months) as recommended, one new food at a time
    \item Continue breastfeeding for as long as possible
    \item Read all food labels and learn the names for milk-derived ingredients
    \item Inform childcare, school, and family members about the allergy
    \item Carry an allergy action plan and, if prescribed, an adrenaline autoinjector
    \item See a dietitian for advice on nutritionally complete alternatives
    \item Do not reintroduce milk without medical advice
    \item Keep regular allergy reviews to check whether the allergy has been outgrown
\end{itemize}

\section*{Sources and Further Reading}
The following reputable sources provide further information for healthcare students and patients seeking reliable, current advice about cow's milk allergy.

\begin{itemize}
    \item Australasian Society of Clinical Immunology and Allergy (ASCIA). \textit{Cow's milk allergy.} https://www.allergy.org.au/
    \item Healthdirect Australia. \textit{Cow's milk allergy in babies and children.} https://www.healthdirect.gov.au/
    \item Allergy \& Anaphylaxis Australia. \textit{Patient information on milk allergy.} https://allergyfacts.org.au/
    \item NHS. \textit{Cow's milk allergy.} https://www.nhs.uk/conditions/cows-milk-allergy/
    \item Mayo Clinic. \textit{Milk allergy: overview.} https://www.mayoclinic.org/diseases-conditions/milk-allergy/
    \item Raising Children Network. \textit{Food allergy in children.} https://raisingchildren.net.au/
\end{itemize}
